% !TEX root = Lecture_Koszul.tex
\clearpage
\section{Koszul algebras}\label{sec:koszul-complex}

Throughout, $A$ is a locally finite-dimensional non-negatively graded algebra with $S=A_0$ finite-dimensional semisimple, as in Subsection~\ref{subsec:graded-ext}.
By Theorem~\ref{thm:minimal-resolution} the module $S=A/A_+$ has a minimal graded projective resolution $P_\bullet\to S$, unique up to isomorphism, and by Theorem~\ref{thm:ext-minimal} and Proposition~\ref{prop:tor-top} it records
\[
\Ext_A^h(S,S)_j\cong\Hom_S\bigl((\grtop P_h)_{-j},S\bigr),
\qquad
\Tor_h^A(S,S)\cong\grtop P_h.
\]
A Koszul algebra is one for which these groups sit as close to the diagonal as possible.

\subsection{Koszul algebras}

\begin{definition}\label{def:linear-resolution}
A graded projective resolution $P_\bullet\to M$ with every $P_h$ bounded below and lfd is \defn{linear} if $P_h$ is generated in degree $h$, that is, $P_h=(P_h)_hA$, for every $h\geq0$.
\end{definition}

A linear resolution is automatically minimal: $(P_{h-1})_h=(P_{h-1})_{h-1}A_1\subseteq P_{h-1}A_+$, and $d_h$ maps $P_h=(P_h)_hA$ into $(P_{h-1})_hA\subseteq P_{h-1}A_+$, so \eqref{eq:graded-minimal-criterion} holds.
Hence minimal resolutions are not needed to define Koszul algebras, only to compute with them.

\begin{definition}\label{def:koszul-algebra}
The algebra $A$ is \defn{Koszul} if the graded right module $S=A/A_+$ has a linear graded projective resolution.
\end{definition}

The generators in homological degree $h$ must then all have grading degree $h$.
Since $S$ can have several simple modules, the definition requires a linear resolution of their direct sum, not just of one chosen simple module.
We do not assume that $A$ is quadratic: this follows from the definition (Corollary~\ref{cor:koszul-quadratic}).

\begin{example}\label{ex:koszul-nonconnected}
The definition also applies when $S$ has several simple modules.
For $A=\Bbbk A_2$, the resolution in Example~\ref{ex:minimal-A2} is linear for $L_1$, and $L_2=e_{22}A$ is projective in degree $0$.
Taking their direct sum gives a linear resolution of $S=L_1\oplus L_2$; hence $\Bbbk A_2$ is Koszul.
\end{example}

\begin{example}
The exterior algebra $E=\Bbbk[\epsilon]/(\epsilon^2)$ is Koszul: Example~\ref{ex:minimal-exterior} gives a resolution of $\Bbbk$ with $P_h=E(-h)$, which is linear although infinite.
\end{example}

For the polynomial algebra a resolution is constructed in Subsection~\ref{subsec:koszul-complex}.

\begin{lemma}\label{lem:resolution-degrees}
The terms of the minimal resolution of $S$ satisfy $(P_h)_i=0$ for $i<h$.
Equivalently, $\Ext_A^h(S,S)_j=0$ for $j>-h$.
\end{lemma}

\begin{proof}
For $h=0$, $P_0=A$.
If $P_h$ lives in degrees $\geq h$, then $\Im d_{h+1}\subseteq P_hA_+$ lives in degrees $\geq h+1$.
So does $\grtop P_{h+1}\cong\grtop\Im d_{h+1}$ (Lemma~\ref{lem:minimal-tops}), hence so does $P_{h+1}$ by \eqref{eq:graded-cover-terms}.
\end{proof}

Thus $A$ is Koszul exactly when the minimal resolution has no generators in degrees $>h$ in homological degree $h$.

\subsection{Characterisation by Ext}

A graded module $M$ is \defn{pure of degree $m$} if $M=M_m$.
Such a module satisfies $MA_+=0$, so it is a right $S$-module, hence semisimple.

\begin{proposition}\label{prop:koszul-ext}
The following are equivalent.
\begin{enumerate}
\item $A$ is Koszul.
\item $\Tor_h^A(S,S)_j=0$ for all $j\neq h$.
\item $\Ext_A^h(S,S)_j=0$ for all $j\neq-h$.
\item For all pure lfd modules $M,N$ of degrees $m,n$, we have $\Ext_A^h(M,N)_j=0$ unless $j=n-m-h$.
\end{enumerate}
\end{proposition}

\begin{proof}
Let $P_\bullet\to S$ be the minimal resolution.
By Proposition~\ref{prop:tor-top}, $\Tor_h^A(S,S)\cong\grtop P_h$.
By Theorem~\ref{thm:ext-minimal}, $\Ext_A^h(S,S)_j\cong\Hom_S((\grtop P_h)_{-j},S)$, which is non-zero if and only if $(\grtop P_h)_{-j}\neq0$, since every simple $S$-module is a summand of $S$.
So (2) and (3) both say that $\grtop P_h$ is concentrated in degree $h$.

(1)$\Rightarrow$(2).
A linear resolution is minimal, hence isomorphic to $P_\bullet$ by Theorem~\ref{thm:minimal-resolution}, and $\grtop P_h=(P_h)_h$.

(2)$\Rightarrow$(1).
By \eqref{eq:graded-cover-terms}, $P_h\cong((\grtop P_h)_h\otimes_SA)(-h)$ is generated in degree $h$.

(1)$\Rightarrow$(4).
Since $M$ is pure and lfd, $M$ is a finite direct sum of modules $L_y(-m)$, each a summand of $S(-m)$, and $\Ext$ is additive in $M$.
Shifting $P_\bullet$ gives $\Ext_A^h(S(-m),N)_j=\Ext_A^h(S,N)_{j+m}$.
As $NA_+=0$, Theorem~\ref{thm:ext-minimal} gives $\Ext_A^h(S,N)_{j'}\cong\Hom_S^\Z(\grtop P_h,N(j'))$.
Here $\grtop P_h$ lives in degree $h$ and $N(j')$ in degree $n-j'$, so this vanishes unless $j'=n-h$, that is, $j=n-m-h$.

(4)$\Rightarrow$(3) is the case $M=N=S$.
\end{proof}

Thus a Koszul algebra is a positively graded algebra that is as close to semisimple as possible: Ext between pure modules is concentrated in one internal degree for each homological degree.
Lemma~\ref{lem:resolution-degrees} shows that for an arbitrary $A$ the groups in (4) already vanish for $j>n-m-h$.

\subsection{Koszul algebras are quadratic}

The groups $\Ext_A^1(S,S)$ and $\Ext_A^2(S,S)$ detect generators and relations, by Proposition~\ref{prop:resolution-first-terms}.

\begin{proposition}\label{prop:koszul-generation}
\begin{enumerate}
\item $A$ is generated in degree $1$ if and only if $\Ext_A^1(S,S)_j=0$ for $j\neq-1$.
\item Suppose that $A$ is generated in degree $1$, and let $I=\Ker(T_S(A_1)\onto A)$.
Then $A$ is quadratic if and only if $\Ext_A^2(S,S)_j=0$ for $j\neq-2$; in this case $A=T_S(A_1)/(I_2)$.
\end{enumerate}
\end{proposition}

\begin{proof}
(1) By Proposition~\ref{prop:resolution-first-terms}(2), $\grtop P_1\cong A_+/A_+^2$, so the vanishing says $A_i=(A_+^2)_i=\sum_{0<p<i}A_pA_{i-p}$ for $i\geq2$.
By induction on $i$ this is equivalent to $A_i=A_1^i$, that is, to surjectivity of all $\mu_i$.

(2) Put $V=A_1$.
By Proposition~\ref{prop:resolution-first-terms}(3), the vanishing says $I_i=V\otimes_SI_{i-1}+I_{i-1}\otimes_SV$ for all $i\geq3$.
By induction on $i$ this gives $I_i\subseteq(I_2)$, so $I=(I_2)$ and $A$ is quadratic.
Conversely, if $A=T_S(V)/(R)$ then $I_2=R$, and Proposition~\ref{prop:resolution-first-terms}(3) shows that $\grtop P_2$ is concentrated in degree $2$.
\end{proof}

\begin{corollary}\label{cor:koszul-quadratic}
Every Koszul algebra is quadratic: $A\cong T_S(A_1)/(R)$ with $R=I_2\subseteq A_1\otimes_SA_1$.
\end{corollary}

\begin{proof}
By Proposition~\ref{prop:koszul-ext}(3), $\Ext_A^1(S,S)$ and $\Ext_A^2(S,S)$ lie on the diagonal.
Apply Proposition~\ref{prop:koszul-generation}.
\end{proof}

The converse fails: not every quadratic algebra is Koszul (Exercise~\ref{ex:quadratic-non-koszul}), so Koszulness is a genuine condition on the $\Ext$ groups in homological degrees $\geq3$.

\subsection{The opposite algebra}

Everything above used right modules, while the definition of $\Tor$ treats both sides of $S$ symmetrically.

\begin{lemma}\label{lem:tor-balanced}
Let $M$ be a graded right and $N$ a graded left $A$-module, both bounded below and lfd, with graded projective resolutions $P_\bullet\to M$ and $Q_\bullet\to N$.
Then $H_h(P_\bullet\otimes_AN)\cong H_h(M\otimes_AQ_\bullet)$ as graded spaces, for all $h\geq0$.
\end{lemma}

\begin{supplementary}
\begin{proof}
Let $C$ be the total complex of the double complex $P_p\otimes_AQ_q$.
Projective modules are flat, so each column $P_p\otimes_AQ_\bullet\to P_p\otimes_AN$ and each row $P_\bullet\otimes_AQ_q\to M\otimes_AQ_q$ is exact.
By the acyclic assembly lemma, the augmentations $C\to P_\bullet\otimes_AN$ and $C\to M\otimes_AQ_\bullet$ are quasi-isomorphisms; they preserve the grading.
\end{proof}
\end{supplementary}

\begin{proposition}\label{prop:opposite-koszul}
$A$ is Koszul if and only if the opposite algebra $A^{\mathrm{op}}$ is Koszul.
\end{proposition}

\begin{proof}
The algebra $A^{\mathrm{op}}$ satisfies the standing hypotheses with $(A^{\mathrm{op}})_0=S^{\mathrm{op}}$, and right $A^{\mathrm{op}}$-modules are left $A$-modules.
By Proposition~\ref{prop:koszul-ext}(2), $A^{\mathrm{op}}$ is Koszul if and only if $\Tor_h^{A^{\mathrm{op}}}(S,S)_j=0$ for $j\neq h$.
Computing with a projective resolution of the left module $S$, this is $H_h(S\otimes_AQ_\bullet)$, which is $\Tor_h^A(S,S)$ by Lemma~\ref{lem:tor-balanced}.
So the condition is the same as in Proposition~\ref{prop:koszul-ext}(2) for $A$.
\end{proof}

Consequently the linear resolution of the right module $S$ exists if and only if one exists for the left module $S$, so ``right Koszul'' and ``left Koszul'' agree.

\subsection{The Koszul complex of a polynomial algebra}\label{subsec:koszul-complex}

The examples so far were small.
To show that polynomial algebras in several variables are Koszul we write down an explicit resolution of $\Bbbk$.

Let $V$ be an $n$-dimensional $\Bbbk$-vector space, choose a basis $x_1,\ldots,x_n$, and let $A=\SymAlg(V)=\Bbbk[x_1,\ldots,x_n]$, with each $x_j$ in degree $1$.
Let $\pi:A\to\Bbbk$ be projection onto degree $0$, so $\Bbbk=A/A_+$.
\phantomsection\label{not:koszul-terms}
For $0\leq h\leq n$, put
\[
K_h:=\bigwedge^h V\otimes_\Bbbk A(-h),
\qquad K_0=A.
\]
The right $A$-action is on the second factor.
The factor $\bigwedge^hV$ is placed in grading degree $0$, so the generators of $K_h$ have grading degree $h$.

\begin{definition}\label{def:koszul-complex}
For $1\leq h\leq n$, define the \defn{Koszul differential}
\begin{equation}\label{eq:koszul-differential}
\begin{aligned}
d_h:K_h&\longrightarrow K_{h-1},\\
(v_1\wedge\cdots\wedge v_h)\otimes a
&\longmapsto
\sum_{r=1}^h(-1)^{r-1}
(v_1\wedge\cdots\wedge\widehat{v_r}\wedge\cdots\wedge v_h)
\otimes v_ra .
\end{aligned}
\end{equation}
The formula is multilinear and alternating in $v_1,\ldots,v_h$, so it defines a map on $\bigwedge^hV$.
The sequence $K_\bullet\to\Bbbk$ consists of these maps followed by $\pi:K_0=A\to\Bbbk$.
\end{definition}

\begin{proposition}\label{prop:koszul-complex-resolution}
The sequence
\[
0\longrightarrow K_n\xrightarrow{d_n}\cdots\xrightarrow{d_2}K_1
\xrightarrow{d_1}K_0\xrightarrow{\pi}\Bbbk\longrightarrow0
\]
is an exact minimal graded free resolution of $\Bbbk$ over $A$.
\end{proposition}

\begin{proof}
Each $d_h$ preserves degree and is right $A$-linear.
For $h\geq2$, compute $d_{h-1}d_h$ on a basis wedge.
Every term is indexed by two distinct positions $r<s$ deleted from that wedge.
Deleting position $r$ and then $s$ gives coefficient $(-1)^{r+s-3}x_{j_s}x_{j_r}$; deleting $s$ and then $r$ gives coefficient $(-1)^{r+s-2}x_{j_r}x_{j_s}$.
Since $x_{j_s}x_{j_r}=x_{j_r}x_{j_s}$, these two terms sum to zero. Hence $d_{h-1}d_h=0$; also $\pi d_1=0$ because every term of $d_1$ is divisible by a variable.

To prove exactness, give $A$ and $\bigwedge V$ the multigrading with $\deg x_j$ equal to the $j$th standard basis vector.
For $\alpha\in\mathbb Z_{\geq0}^n$, let $\operatorname{supp}(\alpha)=\{j\mid\alpha_j>0\}$.
The degree $\alpha$ part of $K_h$ has basis
\[
b_J:=(x_{j_1}\wedge\cdots\wedge x_{j_h})\otimes x^{\alpha-\mathbf1_J},
\qquad
J=\{j_1<\cdots<j_h\}\subseteq\operatorname{supp}(\alpha),
\]
where $\mathbf1_J$ has entry $1$ at indices in $J$ and $0$ elsewhere.
The differential is $d_h(b_J)=\sum_{r=1}^h(-1)^{r-1}b_{J\setminus\{j_r\}}$.

If $\alpha\neq0$, let $q$ be the least index in $\operatorname{supp}(\alpha)$.
For each $h\geq0$, define $s_{\alpha,h}:K_{h,\alpha}\to K_{h+1,\alpha}$ on basis vectors by
\[
s_{\alpha,h}(b_J)=
\begin{cases}
b_{\{q\}\cup J},&q\notin J,\\
0,&q\in J,
\end{cases}
\]
inserting $q$ first in the ordered wedge.
For $q\notin J$ and $h\geq1$, direct substitution gives
\[
\begin{aligned}
d_{h+1}s_{\alpha,h}(b_J)
&=b_J-\sum_{r=1}^h(-1)^{r-1}b_{\{q\}\cup(J\setminus\{j_r\})},\\
s_{\alpha,h-1}d_h(b_J)
&=\sum_{r=1}^h(-1)^{r-1}b_{\{q\}\cup(J\setminus\{j_r\})}.
\end{aligned}
\]
For $h=0$, $J=\varnothing$ and $d_1s_{\alpha,0}(b_\varnothing)=b_\varnothing$.
For $q\in J$, the term of $d_h(b_J)$ deleting $q$ is the only term on which $s_{\alpha,h-1}$ is nonzero, and it gives $s_{\alpha,h-1}d_h(b_J)=b_J$ while $d_{h+1}s_{\alpha,h}(b_J)=0$.
Thus $d_{h+1}s_{\alpha,h}+s_{\alpha,h-1}d_h=\id_{K_{h,\alpha}}$ for every $h\geq1$ in positive multidegree; the case $h=0$ was checked separately.
All positive multidegrees are therefore exact.
In multidegree $0$, only $K_0\cong\Bbbk$ is nonzero and $\pi$ is the identity onto $\Bbbk$.
This proves exactness of the displayed sequence.

Finally, each $K_h$ is a finite direct sum of copies of $A(-h)$, and \eqref{eq:koszul-differential} has image in $K_{h-1}A_+$.
Also $\Ker\pi=AA_+$.
Thus $\Im d_h\subseteq K_{h-1}A_+$ for all $h\geq1$ and $\Ker\pi\subseteq K_0A_+$, so the resolution is minimal by \eqref{eq:graded-minimal-criterion}.
\end{proof}

\begin{proposition}\label{prop:polynomial-koszul}
For every finite-dimensional vector space $V$, the symmetric algebra $\SymAlg(V)$ is Koszul.
\end{proposition}

\begin{proof}
Proposition~\ref{prop:koszul-complex-resolution} gives a minimal resolution with $K_h=\bigwedge^hV\otimes_\Bbbk A(-h)$.
Thus $K_h$ is generated in degree $h$ for every $h$.
\end{proof}

Proposition~\ref{prop:tor-top} gives $\Tor_h^A(\Bbbk,\Bbbk)\cong\grtop K_h=\bigwedge^hV$, so $\Ext_A^h(\Bbbk,\Bbbk)_{-h}\cong D(\bigwedge^hV)\cong\bigwedge^hD(V)$.
This is the degree $h$ part of $\SymAlg(V)^!\cong\bigwedge D(V)$ from Proposition~\ref{prop:symmetric-exterior-dual}: the quadratic dual computed combinatorially in Section~\ref{sec:quadratic-dual} appears here as an $\Ext$ algebra.


\subsection{Exercises}

\begin{exercise}\label{ex:hilbert-polynomial}
Let $A=\SymAlg(V)=\Bbbk[x_1,\ldots,x_n]$ and $K_\bullet\to\Bbbk$ be the Koszul resolution of Proposition~\ref{prop:koszul-complex-resolution}.
In each degree, the alternating sum of the dimensions in an exact sequence of finite-dimensional vector spaces is zero.
Use this to show
\[
(1-t)^n\,h_A(t)=1,
\]
and deduce the Hilbert series of $\Bbbk[x_1,\ldots,x_n]$.
\end{exercise}

\begin{exercise}
Let $S=\bigoplus_{y=1}^n\Bbbk e_y$, let $V$ be a finite-dimensional bimodule in degree $1$, and let $A=T_S(V)/(V\otimes_SV)=S\oplus V$.
Prove that the minimal graded projective resolution of $S$ has $P_h\cong(V^{\otimes_Sh}\otimes_SA)(-h)$ for all $h\geq0$, and deduce that $A$ is Koszul.

\emph{Hint.} Show that $\Ker p_0=V$ and that $VA_+=0$, so each kernel is killed by $A_+$.
\end{exercise}

\begin{exercise}
Let $A=\Bbbk[x]/(x^3)$ with $\deg x=1$.
Find the minimal graded projective resolution of $\Bbbk$, and show that $P_h$ is generated in degree $\lfloor 3h/2\rfloor$ for all $h\geq0$.
Deduce that $A$ is not Koszul.
Is $A$ quadratic?
\end{exercise}

\begin{exercise}\label{ex:quadratic-non-koszul}
Let $A=\Bbbk\langle x,y\rangle/(yx,\;x^2+y^2)$ with $\deg x=\deg y=1$, a quadratic algebra with $S=\Bbbk$.
\begin{enumerate}
\item Show that $x^3=0$ and $y^3=0$, and deduce that $A_n=0$ for $n\geq3$ and $h_A(t)=1+2t+2t^2$.
\item Suppose that $A$ is Koszul, with linear resolution $P_h=(T_h\otimes_\Bbbk A)(-h)$ and $b_h=\dim_\Bbbk T_h$.
By taking the alternating sum of Hilbert series in the exact sequence $\cdots\to P_1\to P_0\to\Bbbk\to0$ as in Exercise~\ref{ex:hilbert-polynomial}, show that $\bigl(\sum_{h\geq0}(-1)^hb_ht^h\bigr)\,h_A(t)=1$.
\item Expand $1/h_A(t)$ to degree $4$ and obtain a contradiction.
Conclude that $A$ is quadratic but not Koszul.
\end{enumerate}
\emph{Hint for (1).} $x^3=x\cdot x^2=-xy^2$ and $x^3=x^2\cdot x=-y^2x=-y(yx)$, and similarly for the other words of length $3$.
\end{exercise}

\begin{hidden}
\emph{Solution to the last exercise.}
(1) We have $x^3=x\cdot x^2=-xy^2$ and $x^3=x^2\cdot x=-y^2x=-y(yx)=0$, so $x^3=0$ and $xy^2=0$.
Likewise $y^3=y\cdot y^2=-yx^2=-(yx)x=0$.
The remaining words of length $3$ contain $yx$ or are $xxy=-y^3=0$ and $xyy=-x^3=0$, hence vanish.
So $A_3=0$, and $A_n=0$ for $n\geq3$ since $A$ is generated in degree $1$.
Also $\dim A_1=2$ and $\dim A_2=4-\dim R=2$, so $h_A(t)=1+2t+2t^2$.

(2) Taking the degree $i$ component of $\cdots\to P_1\to P_0\to\Bbbk\to0$ and the alternating sum of dimensions gives $\sum_h(-1)^h\dim(P_h)_i=\dim\Bbbk_i$.
In generating functions, $h_{P_h}(t)=b_ht^hh_A(t)$ and $h_\Bbbk(t)=1$, so $\bigl(\sum_h(-1)^hb_ht^h\bigr)h_A(t)=1$.

(3) Put $u(t)=\sum_h(-1)^hb_ht^h=1/(1+2t+2t^2)$.
Then $u_0=1$, $u_1=-2$, and $u_h=-2u_{h-1}-2u_{h-2}$ gives $u_2=2$, $u_3=0$, $u_4=-4$.
So $b_4=u_4=-4<0$, which is impossible.
Hence $A$ has no linear resolution of $\Bbbk$, so $A$ is not Koszul, although it is quadratic.
\end{hidden}
